%% kapa: exam=1 round=24 year=2013 period=1 form=A subject=civil_law subject_ko=민법 seq=4 number=4
%% source: Q-Net 감정평가사 기출문제 (cst00309) artlSeq=5073802 file=extracted/1008/5073802/2013년 제24회 감정평가사 1차시험 기출문제/1교시(A형).hwp
%% conversion: kordoc 4.24.51 -> tools/gen1.py (auto); review=proofread-sonnet; flags=none
\begin{question}{4}
법인에 관한 설명으로 옳지 않은 것은? (다툼이 있으면 판례에 의함)
\begin{choices}
  \item 법인 대표자의 행위가 직무에 해당하지 않음을 그 행위 당시 알았던 피해자는 법인의 불법행위책임을 물을 수 없다.
  \item 사단법인의 정관에 정관변경을 금지하는 규정이 있더라도 사원총회의 결의로 그 규정을 변경할 수 있다.
  \item 이사가 여럿인 경우 정관에 다른 규정이 없으면 이사의 과반수로서 법인을 대표하여야 한다.
  \item 각 사원은 원칙적으로 평등한 결의권을 갖지만, 정관으로 달리 정할 수 있다.
  \item 재단법인이 정관에 정한 방법으로 정관을 변경하더라도 주무관청의 허가가 없으면 그 효력이 없다.
\end{choices}
\end{question}
\answer{3}
